% !TEX root =./ECE448_Practice_main.tex
%
% L15 practice set (mPSK). Every question is a step in the L15 link: bits per
% symbol, the Gray map, differential coding by hand, the frequency
% offset the Costas Loop has to remove, and choosing between QPSK and 8PSK.
% Blank copy: \begin{solution}[h] reserves h of work space.
% SOLUTIONS copy (\iskey=1) prints the red worked solutions.

\fullwidth{\textbf{Learning Objective 6.4: } I can implement an mPSK modulator and
demodulator (BPSK, QPSK, 8PSK), recover carrier and symbol timing, and verify
constellation behavior against theory.
\\[4pt]\normalfont\small Show your work. Write a \textbf{Documentation} line at the
end (\textbf{None} if you did not collaborate).}

\begin{questions}

%%%%%%%%%%%%%%%%%%%%%%%%%%%%%%%%%%%%%%%%%
\question \textbf{Bits per symbol.} A telemetry radio sends 25,000 symbols per second
with RRC pulses of roll-off $\alpha = 0.25$. On a circle of radius 1, neighboring points
sit 2.00 apart for BPSK, 1.41 apart for QPSK, and 0.77 apart for 8PSK (Lesson 3).
\begin{parts}

	\begin{minipage}[t]{0.58\linewidth}
		\part Find the occupied bandwidth, and the bit rate with BPSK, QPSK, and 8PSK.
		\begin{solution}[1.2in]
		\eq{ B &= (1 + \alpha) R_s = 1.25 \times 25\ \text{kHz} = 31.25\ \text{kHz} }
		The bandwidth does not depend on $M$. The bit rate is $k R_s$ with $k = \log_2 M$:
		25, 50, and 75 kb/s.
		\end{solution}
	\end{minipage}\hfill%
	\begin{minipage}[t]{0.37\linewidth}
		\raggedleft \vspace{12pt}
		$B$: \ansbox[1.0in]{31.25 kHz}

		\vspace{4pt}
		bits: \ansbox[1.3in]{25, 50, 75 kb/s}
	\end{minipage}

	\begin{minipage}[t]{0.58\linewidth}
		\part The telemetry must carry 60 kb/s. Which of the three modulations can carry it
		at 25,000 symbols per second?
		\begin{solution}[1.0in]
		It needs at least $60/25 = 2.4$ bits per symbol. BPSK carries 1 and QPSK 2, so
		only 8PSK, with 3 bits and 75 kb/s, can carry it.
		\end{solution}
	\end{minipage}\hfill%
	\begin{minipage}[t]{0.37\linewidth}
		\raggedleft \vspace{12pt}
		\ansbox[1.0in]{8PSK}
	\end{minipage}

	\part In words, explain what the radio gives up by moving from QPSK to 8PSK, and
	what you would watch on the receiver's constellation display as the signal weakens.
		\begin{solution}[1.2in]
		8PSK sends 50\% more bits in the same 31.25 kHz, but its points sit 0.77 apart
		instead of 1.41, so the same noise pushes far more symbols across to a neighbor.
		As the signal weakens, the noise clouds grow; the eight 8PSK clouds touch and
		merge into a ring while QPSK's four clouds are still separate, and the text
		starts to break up sooner.
		\end{solution}

\end{parts}

\newpage
%%%%%%%%%%%%%%%%%%%%%%%%%%%%%%%%%%%%%%%%%
\question \textbf{The Gray map.} The L15 transmitter numbers the QPSK points by
\emph{phase index} $i = 0, 1, 2, 3$, counterclockwise from the point at $45^\circ$.
The Gray label of index $i$ is $i \oplus (i \gg 1)$, the index XORed with itself
shifted right one bit.
\begin{parts}

	\begin{minipage}[t]{0.58\linewidth}
		\part Write the 2-bit Gray label of each index. Then find the indices the
		transmitter sends for the bits \texttt{01 11 10 00}.
		\begin{solution}[1.4in]
		\eq{ 0 \oplus 0 = 00, &\quad 1 \oplus 0 = 01 \\
		2 \oplus 1 = 11, &\quad 3 \oplus 1 = 10 }
		The labels are 00, 01, 11, 10. The transmitter's Map sends the index whose label
		matches each pair: \texttt{01} is 1, \texttt{11} is 2, \texttt{10} is 3, and
		\texttt{00} is 0, so it sends 1, 2, 3, 0.
		\end{solution}
	\end{minipage}\hfill%
	\begin{minipage}[t]{0.37\linewidth}
		\raggedleft \vspace{12pt}
		labels: \ansbox[1.3in]{00, 01, 11, 10}

		\vspace{4pt}
		indices: \ansbox[1.3in]{1, 2, 3, 0}
	\end{minipage}

	\begin{minipage}[t]{0.58\linewidth}
		\part Noise moves a symbol sent as index 1 into the region of index 2. How many
		bits are wrong with the Gray labels? How many would be wrong if index $i$
		simply carried the binary number $i$?
		\begin{solution}[1.0in]
		Gray: \texttt{01} becomes \texttt{11}, one bit. Natural binary: \texttt{01}
		becomes \texttt{10}, two bits.
		\end{solution}
	\end{minipage}\hfill%
	\begin{minipage}[t]{0.37\linewidth}
		\raggedleft \vspace{12pt}
		Gray: \ansbox[1.0in]{1}

		\vspace{4pt}
		binary: \ansbox[1.0in]{2}
	\end{minipage}

	\part In words, explain why a wrong 8PSK symbol in a Gray-coded link usually
	corrupts only one of its three bits.
		\begin{solution}[1.0in]
		Almost every symbol error is a slip to one of the two neighbors, because the
		noise rarely reaches farther. With Gray labels a neighbor differs in one of the
		three bits, so each wrong symbol usually corrupts one bit out of three.
		\end{solution}

\end{parts}

\newpage
%%%%%%%%%%%%%%%%%%%%%%%%%%%%%%%%%%%%%%%%%
\question \textbf{Differential coding by hand.} The L15 Differential Encoder sends
$e_n = (e_{n-1} + d_n) \bmod M$, where $d_n$ is the phase index from the Gray map and
$e_{n-1}$ starts at 0. The decoder outputs $(r_n - r_{n-1}) \bmod M$, where $r_n$ is the
received index and $r_{n-1}$ also starts at 0. The link uses QPSK, and the data indices
are $d = 2, 1, 1, 3, 0$.
\begin{parts}

	\begin{minipage}[t]{0.58\linewidth}
		\part Find the five indices the encoder sends.
		\begin{solution}[1.0in]
		\eq{ e_0 &= 0 + 2 = 2, \quad e_1 = 2 + 1 = 3 \\
		e_2 &= (3 + 1) \bmod 4 = 0 \\
		e_3 &= 0 + 3 = 3, \quad e_4 = 3 + 0 = 3 }
		\end{solution}
	\end{minipage}\hfill%
	\begin{minipage}[t]{0.37\linewidth}
		\raggedleft \vspace{12pt}
		\ansbox[1.3in]{2, 3, 0, 3, 3}
	\end{minipage}

	\begin{minipage}[t]{0.58\linewidth}
		\part The receiver's Costas Loop locks $180^\circ$ off, so every received index
		is 2 more than the one sent, mod 4. Decode the five received indices and compare
		them with $d$.
		\begin{solution}[1.4in]
		Received: $r = 0, 1, 2, 1, 1$. Decoded, with $r_{-1} = 0$:
		\eq{ 0 - 0 = 0,\ 1 - 0 = 1,\ 2 - 1 = 1,\ 1 - 2 \equiv 3,\ 1 - 1 = 0 }
		The output is 0, 1, 1, 3, 0. Every symbol after the first is right; the first is
		wrong only because the decoder's starting reference was not rotated.
		\end{solution}
	\end{minipage}\hfill%
	\begin{minipage}[t]{0.37\linewidth}
		\raggedleft \vspace{12pt}
		\ansbox[1.3in]{0, 1, 1, 3, 0}
	\end{minipage}

	\begin{minipage}[t]{0.58\linewidth}
		\part Now the loop locks correctly, but noise turns the third received index from
		0 into 1. Decode the five indices and count the wrong symbols.
		\begin{solution}[1.2in]
		Received: $r = 2, 3, 1, 3, 3$.
		\eq{ 2,\ 3 - 2 = 1,\ 1 - 3 \equiv 2,\ 3 - 1 = 2,\ 3 - 3 = 0 }
		The output is 2, 1, 2, 2, 0: the third and fourth symbols are wrong, two errors
		from one wrong decision.
		\end{solution}
	\end{minipage}\hfill%
	\begin{minipage}[t]{0.37\linewidth}
		\raggedleft \vspace{12pt}
		\ansbox[1.0in]{2 wrong}
	\end{minipage}

	\part In words, use parts (b) and (c) to explain what differential coding does for
	a link and what it does to the errors in a noisy link.
		\begin{solution}[1.0in]
		It makes the link work at any of the $M$ lock positions, because a rotation adds
		the same amount to every index and the differences do not change. In exchange,
		every wrong decision corrupts two differences, the one into it and the one out
		of it, so a noisy link makes about twice as many errors.
		\end{solution}

\end{parts}

\newpage
%%%%%%%%%%%%%%%%%%%%%%%%%%%%%%%%%%%%%%%%%
\question \textbf{The frequency offset.} A frequency offset $\Delta f$ turns the
constellation by $360^\circ \times \Delta f / R_s$ every symbol. Lesson 15 measured the
pull-in range of the QPSK Costas Loop at about 375 Hz for $R_s = 4{,}000$ symbols per
second and the default noise amplitude of 0.5. The loop works once per symbol, so with the same
loop settings its pull-in range scales with the symbol rate.
\begin{parts}

	\begin{minipage}[t]{0.58\linewidth}
		\part A 4 kbaud QPSK link has a 75 Hz offset. Find the rotation per symbol and the
		number of symbols in one full turn of the constellation.
		\begin{solution}[1.0in]
		\eq{ 360^\circ \times \dfrac{75}{4000} &= 6.75^\circ \\
		\dfrac{360^\circ}{6.75^\circ} &= 53.3\ \text{symbols} }
		\end{solution}
	\end{minipage}\hfill%
	\begin{minipage}[t]{0.37\linewidth}
		\raggedleft \vspace{12pt}
		\ansbox[1.0in]{$6.75^\circ$}

		\vspace{4pt}
		\ansbox[1.0in]{53.3}
	\end{minipage}

	\begin{minipage}[t]{0.58\linewidth}
		\part Two radios use oscillators accurate to 2.5 ppm and transmit at 2.4 GHz. Find
		the largest frequency offset between them. Can a 4 kbaud QPSK Costas Loop lock
		to it? Can a 1 Mbaud one?
		\begin{solution}[1.4in]
		Each can be off by $2.5 \times 10^{-6} \times 2.4\ \text{GHz} = 6$ kHz, in
		opposite directions:
		\eq{ \Delta f_\text{max} &= 2 \times 6\ \text{kHz} = 12\ \text{kHz} }
		At 4 kbaud the loop pulls in from about 375 Hz, 32 times too little. At 1 Mbaud
		the range scales to $375 \times 250 \approx 94$ kHz, which covers 12 kHz.
		\end{solution}
	\end{minipage}\hfill%
	\begin{minipage}[t]{0.37\linewidth}
		\raggedleft \vspace{12pt}
		\ansbox[1.0in]{12 kHz}

		\vspace{4pt}
		4 kbaud: \ansbox[0.8in]{no}

		\vspace{4pt}
		1 Mbaud: \ansbox[0.8in]{yes}
	\end{minipage}

	\part In words, explain why an 8PSK Costas Loop pulls in from a smaller offset than
	a BPSK loop with the same bandwidth, and what the 4 kbaud receiver in part (b)
	needs before its Costas Loop.
		\begin{solution}[1.2in]
		The phase detector cannot tell one constellation point from another, so it can
		only measure a phase error up to half the spacing, $180^\circ/M$: $90^\circ$ for
		BPSK but $22.5^\circ$ for 8PSK. A spinning constellation leaves the detector's
		range faster when $M$ is larger, so the 8PSK loop has less time to correct and
		locks only from a smaller offset. The 4 kbaud receiver needs a coarse frequency
		correction first, such as an FLL Band-Edge block or a frequency measured from a
		known preamble, to bring the offset inside the Costas Loop's range.
		\end{solution}

\end{parts}

\newpage
%%%%%%%%%%%%%%%%%%%%%%%%%%%%%%%%%%%%%%%%%
\question \textbf{QPSK or 8PSK?} A link must carry 12 kb/s in a channel 6.25 kHz wide,
with RRC pulses of roll-off $\alpha = 0.35$ and differential coding. Lesson 15's Noise frame
showed QPSK and 8PSK at the same noise amplitude, 0.8.
\begin{parts}

	\begin{minipage}[t]{0.58\linewidth}
		\part Find the symbol rate and the occupied bandwidth with QPSK and with 8PSK.
		Which fits the channel?
		\begin{solution}[1.4in]
		\eq{ \text{QPSK: } R_s &= 12/2 = 6\ \text{kbaud}, \quad B = 1.35 \times 6 = 8.1\ \text{kHz} \\
		\text{8PSK: } R_s &= 12/3 = 4\ \text{kbaud}, \quad B = 1.35 \times 4 = 5.4\ \text{kHz} }
		Only 8PSK fits in 6.25 kHz.
		\end{solution}
	\end{minipage}\hfill%
	\begin{minipage}[t]{0.37\linewidth}
		\raggedleft \vspace{12pt}
		QPSK: \ansbox[1.0in]{8.1 kHz}

		\vspace{4pt}
		8PSK: \ansbox[1.0in]{5.4 kHz}
	\end{minipage}

	\part Describe what the receiver's constellation and Received label showed for each
	modulation at a noise amplitude of 0.8, and what that means for the 8PSK link.
		\begin{solution}[1.2in]
		QPSK's four clouds nearly touched but stayed separate, and most lines arrived.
		8PSK's eight clouds merged into a ring, and almost no line arrived intact. The 8PSK
		link fits the channel but needs a quieter channel, a stronger signal, than QPSK
		would at the same noise.
		\end{solution}

	\part In words, state which modulation you would choose and what you would check on
	the bench before committing to it.
		\begin{solution}[1.2in]
		8PSK, because QPSK does not fit the channel at 12 kb/s. On the bench, run the link at
		the weakest signal you expect and check that the eight clouds stay separate on the
		constellation display and the text arrives intact, with some room to spare. If they
		merge, the link needs more power or a slower bit rate.
		\end{solution}

\end{parts}

\end{questions}
